%~~~~~~~~~~~~~~~~~~~~~~~~~~~~~%
%  dernière                   %
%  modification = 26/04/2023  %
%~~~~~~~~~~~~~~~~~~~~~~~~~~~~~%
%      Colas Bardavid


% --------------------------------------------------------- %
%                       Physique
% --------------------------------------------------------- %

\DeclareMathOperator{\incertitude}{\textsf{\textup{u}}}
\def\zScore{\textsf{\textup{z}}}

% --------------------------------------------------------- %
%                           Listes
% --------------------------------------------------------- %

\newlist{listeBullet}{itemize}{1}
\setlist[listeBullet]{label = \textbullet}

\newlist{listeEtoile}{itemize}{1}
\setlist[listeEtoile]{label = $\circledast$}


% --------------------------------------------------------- %
%              Diverses macros pour les maths
% --------------------------------------------------------- %

\def\Q{\mathbb{Q}}
\def\R{\mathbb{R}}
\def\Z{\mathbb{Z}}
\def\N{\mathbb{N}}
\def\C{\mathbb{C}}
\def\K{\mathbb{K}}
\def\U{\mathbb{U}}

\renewcommand{\d}[1]{\ensuremath{\operatorname{d}\!{#1}}} % le d des intégrales

\def\iC{\textrm{i}} % le i des nombres complexes
\def\jC{\textrm{j}} % le j des nombres complexes

\def\eEuler{\textrm{e}} % le e de Euler

\newcommand{\etc}{\emph{etc.}}

\let\Im\relax
\let\Re\relax
\DeclareMathOperator{\Im}{\mathsf{Im}}
\DeclareMathOperator{\Re}{\mathsf{Re}}

\def\longto{\mathop{\longrightarrow}}

\newcommand{\cvg}[1]{\xrightarrow[#1]{}}


% --------------------------------------------------------- %
%                  Lettres pour les points
% --------------------------------------------------------- %

\newcommand{\macroPourLesPoints}[1]{\text{$\mathrm{#1}$} }
\def\ptA{\macroPourLesPoints{A}}
\def\ptB{\macroPourLesPoints{B}}
\def\ptC{\macroPourLesPoints{C}}
\def\ptD{\macroPourLesPoints{D}}
\def\ptE{\macroPourLesPoints{E}}
\def\ptF{\macroPourLesPoints{F}}
\def\ptG{\macroPourLesPoints{G}}
\def\ptH{\macroPourLesPoints{H}}
\def\ptI{\macroPourLesPoints{I}}
\def\ptJ{\macroPourLesPoints{J}}
\def\ptK{\macroPourLesPoints{K}}
\def\ptL{\macroPourLesPoints{L}}
\def\ptM{\macroPourLesPoints{M}}
\def\ptN{\macroPourLesPoints{N}}
\def\ptO{\macroPourLesPoints{O}}
\def\ptP{\macroPourLesPoints{P}}
\def\ptQ{\macroPourLesPoints{Q}}
\def\ptR{\macroPourLesPoints{R}}
\def\ptS{\macroPourLesPoints{S}}
\def\ptT{\macroPourLesPoints{T}}
\def\ptU{\macroPourLesPoints{U}}
\def\ptV{\macroPourLesPoints{V}}
\def\ptW{\macroPourLesPoints{W}}
\def\ptX{\macroPourLesPoints{X}}
\def\ptY{\macroPourLesPoints{Y}}
\def\ptZ{\macroPourLesPoints{Z}}
\def\ptOmega{\macroPourLesPoints{\Omega}}


% --------------------------------------------------------- %
%          Commandes pour dessiner les fonctions
% --------------------------------------------------------- %
% Exemples :
% \fonction{\R}{\R}{x}{x^2}				fonction R->R, x |-> x^2
% \fonctionAccolade{\R}{\R}{x}{x^2}		fonction R->R, x |-> x^2 avec une accolade
%
% Warning: xypic changed to xy
\RequirePackage[all]{xy}
\makeatother
\newlength\longueurdudiagramme
\newcommand{\fonction}[4]
{%
\settowidth{\longueurdudiagramme}
{%
$\xymatrix@R=0mm{
{#1} \ar[r] & {#2} \\
{#3} \ar@{|->}[r] & {#4}}$
}
\begin{minipage}{\longueurdudiagramme}
$\xymatrix@R=0mm{
{#1} \ar[r] & {#2} \\
{#3} \ar@{|->}[r] & {#4}}$
\end{minipage}
}
\makeatletter

\newcommand{\fonctionAccolade}[4]
{\left\{
\fonction{#1}{#2}{#3}{#4}
\right.}



% --------------------------------------------------------- %
%                       Guillemets
% --------------------------------------------------------- %

\newcommand{\glm}[1]{\enquote{#1}}



% --------------------------------------------------------- %
%                  Signe d'équivalence
% --------------------------------------------------------- %

\def\ssi{\iff}
\def\quadSSI{\quad\ssi\quad}
\def\qquadSSI{\qquad\ssi\qquad}
\def\spaceSSI{\ \ssi\ }



% --------------------------------------------------------- %
%                      Valeur absolue
% --------------------------------------------------------- %

\DeclarePairedDelimiter{\myVA}{\lvert}{\rvert}
\newcommand{\va}[1]{\myVA*{#1}}
\newcommand{\vaN}[1]{\myVA[\normalsize]{#1}}
\newcommand{\vaB}[1]{\myVA[\big]{#1}}
\newcommand{\vaBB}[1]{\myVA[\Big]{#1}}
\newcommand{\vaBBB}[1]{\myVA[\bigg]{#1}}
\newcommand{\vaBBBB}[1]{\myVA[\Bigg]{#1}}




% --------------------------------------------------------- %
%               Macros perso de Jimmy Roussel
% --------------------------------------------------------- %

% commande \resistanceH{option}{nom}  : résistance de longueur unité, centrée en (0,0).
\newcommand{\resistanceH}[2]{
	\draw[#1,fill=white] (-0.5,3pt)--++(1,0)--++(0,-6pt)--++(-1,0)--cycle;
	\draw[#1] (0,0) node[above=3pt]{#2};
}

% commande \resistanceV{option}{nom}  : résistance de longueur unité, centrée en (0,0).
\newcommand{\resistanceV}[2]{
	\draw[#1,fill=white] (-3pt,0.5)--++(0,-1)--++(6pt,0)--++(0,1)--cycle;
	\draw[#1] (0,0) node[left=3pt]{#2};
}
% \pile{options}{valeur} : pile avec le voltage à gauche
\newcommand{\pileL}[2]%
{
\draw[#1,draw=white,very thick](0,.5)--(0,-.5);
\draw[#1](0,.5)--(0,2pt);
\draw[#1](0,-.5)--(0,-2pt);
\draw[#1] (-10pt,2pt)--++(20pt,0)node[right]{\scriptsize +};
\draw[#1,ultra thick] (-3pt,-2pt)--++(6pt,0);
\draw[#1](-10pt,0)node[left]{ #2};
}
% \pile{options}{valeur} : pile avec le voltage à droite
\newcommand{\pileR}[2]%
{
\draw[#1,draw=white,very thick](0,.5)--(0,-.5);
\draw[#1](0,.5)--(0,2pt);
\draw[#1](0,-.5)--(0,-2pt);
\draw[#1] (-10pt,2pt)node[left]{\scriptsize +}--++(20pt,0);
\draw[#1,ultra thick] (-3pt,-2pt)--++(6pt,0);
\draw[#1](10pt,0)node[right]{ #2};
}
% \voltmetre{options} : voltmètre de longueur unité, centrée en (0,0)
\newcommand{\voltmetre}[1]%
{
\draw[thick,fill=white,#1] (-.5,-.25)rectangle(.5,.75);
\draw[#1] (0,.75)node[above]{Voltmètre};
\draw[#1,fill=lightgray] (0,-.25)++(120:9mm)arc(120:60:9mm)--++(-120:2mm)arc(60:120:7mm)--++(120:2mm);
\draw[#1,thick](0,-.25)++(80:7mm)--++(80:1.5mm);
\draw[#1](-.5,0)--++(10pt,0);
\draw[#1](.5,0)--++(-10pt,0);
\draw[#1,fill=white](-.5cm+10pt,0)circle(2pt);
\draw[#1,fill=white](.5cm-10pt,0)circle(2pt);
\draw[#1](.5cm-12pt,5pt)--++(4pt,0);
\draw[#1](-.5cm+12pt,5pt)--++(-4pt,0);
\draw[#1](-.5cm+10pt,7pt)--++(0,-4pt);
}

% \FEMVL{options}{valeur à gauche}: source de tension verticale
\newcommand{\FEMVL}[2]%
{
\draw[#1,fill=white] (0,0) circle(0.3);
\draw[#1] (0,-0.5)--++(0,1);
\draw[#1,->] (-4mm,-3mm)--(-4mm,3mm)node[midway,left]{\footnotesize #2};
}

% \FEMVR{options}{valeur à droite}: source de tension verticale
\newcommand{\FEMVR}[2]%
{
\draw[#1,fill=white] (0,0) circle(0.3);
\draw[#1] (0,-0.5)--++(0,1);
\draw[#1,->] (4mm,-3mm)--(4mm,3mm)node[midway,right]{\footnotesize #2};
}
